\rok{Фебруар 2024.}{В - Поправни тест првог теста}

Други поправни тест из Алгоритама и структура података

Први практични тест одржан 30.01.2024.

\zadatak{1}{Брисање елемената из реда (Queue Dequeue)}
Написати методу која имплементира стандардан начин за брисање елемената из реда.

\textbf{Додатне напомене за решавање задатка:}
\begin{itemize}
	\item Алгоритам написати у оквиру закоментарисане методе \texttt{dequeue?()}.
	\item Поменута метода налази се у оквиру класе \texttt{Queue} која се налази у придруженом template-у.
	\item Обезбедити код у Main методи за тестирање алгоритма.
\end{itemize}



\zadatak{2}{Симетричност матрице}
Написати алгоритам који ће проверити да ли је матрица симетрична. Нека квадратна матрица је симетрична ако се елементи изнад и испод главне дијагонале пресликавају. Главна дијагонала садржи све елементе на путањи од поља које се налази у горњем левом ћошку, до поља које се налази у доњем десном ћошку.

$$\begin{bmatrix} 0 & 3 & 2 \\ 3 & 0 & 4 \\ 2 & 4 & 0 \end{bmatrix}$$

\textbf{Додатни захтеви за решавање задатка:}
\begin{itemize}
	\item Алгоритам треба да резултује исписом у конзоли у формату: „Matrica (ni)je simetrična“.
	\item Алгоритам писати у оквиру методе:
	\begin{Verbatim}[fontsize=\footnotesize, numbers=left, xleftmargin=10mm]
public static void isMatrixSymmetric(long[,] matrix, int rows, int cols)
	\end{Verbatim}
\end{itemize}

\textbf{Пример:}
\begin{itemize}
	\item \textbf{Улаз 1:} $\begin{bmatrix} 0 & -1 & 2 \\ 3 & 0 & 4 \\ 2 & 4 & 0 \end{bmatrix}$, 3, 3
	\item \textbf{Излаз 1:} „Matrica nije simetrična“
	\item \textbf{Улаз 2:} $\begin{bmatrix} 0 & 3 & 2 \\ 3 & 0 & 4 \\ 2 & 4 & 0 \end{bmatrix}$, 3, 3
	\item \textbf{Излаз 2:} „Matrica je simetrična“
\end{itemize}



\zadatak{3}{Проналажење најдуже подлисте парних бројева}
Написати алгоритам који проналази најдужу подлисту парних елемената. Вредности из пронађене подлисте исписати у конзоли, или исписати да таква подлиста не постоји.

\textbf{Додатне напомене за решавање задатка:}
\begin{itemize}
	\item Захтеван потпис методе:
	\begin{Verbatim}[fontsize=\footnotesize, numbers=left, xleftmargin=10mm]
public void findLargestEvenNumbersSubList()
	\end{Verbatim}
	\item Методу писати у оквиру класе \texttt{LinkedList}.
\end{itemize}

\textbf{Примери:}
\begin{itemize}
	\item \textbf{Улаз:} \texttt{LinkedList: 1 -> 0 -> 6 -> 8 -> 4}
	\item \textbf{Излаз:} \texttt{0 6 8 4}
	\item \textbf{Улаз:} \texttt{LinkedList: 1 -> 5 -> 3}
	\item \textbf{Излаз:} \texttt{Ne postoji podlista parnih brojeva}
\end{itemize}



\sablon{Шаблон поправног за први тест фебруар 2024. група Б}{2023/Februar/GrupaB/AISP2023_PopravniFebruar_Test1_GrupaB.linq}
